\scanpage{44}
\margin{43}
On posera
\[
(3)\qquad D=D_{\widetilde{\unsure{I|J}}}=\operatorname{Aut}(V,\widetilde{\unsure{I|J}}).
\]
\[
(4)\qquad\underline\omega=\underline\omega_{\widetilde{\unsure{I|J}}}
=\Omega(\widetilde{\unsure{I|J}})\simeq\Omega^{\unsure{\vee}}(V)
\]
\margin{ens des orientations du poly. comb. $\widetilde{\unsure{I|J}}$,
ou de l'espace vectoriel réel $V$}

On a dans $D$ le sous-groupe $D^+$ des rotations, qui est le noyau
de l'hom naturel \inserted{(transp. de str.)}
$D\longrightarrow\mathbb Z/2\mathbb Z\simeq\operatorname{Aut}(\omega)$,
d'où suite exacte
\diagram{ (5)\quad\begin{tikzcd}[ampersand replacement=\&, column sep=large]
1\arrow[r,dash]\&D^+\arrow[r]\&D\arrow[r]\&\mathbb Z/2\mathbb Z\arrow[r,dash]\&1
\end{tikzcd}}
$\mathbb Z/2\mathbb Z$ opérant sur $D^+$ par $\zeta\mapsto\zeta^{-1}$.
On \deleted{y a un point} \inserted{a un plongement}
\[
(6)\qquad\omega\subset D^+
\]
et $\underline\omega$ est formé de deux rotations inverses l'une de l'autre,
d'angle $\dfrac{2\pi}{m}$ \deleted{pour l'} \inserted{dans le} sens de
\deleted{l'} \inserted{l'une et l'autre} orientations du $V$.
Ce sont deux générateurs de $D^+$, inverses
\deleted{\unsure{de l'un}\ \unsure{On si}}
\inserted{l'un de l'autre $\omega'=\omega^{-1}$,}
\inserted{le choix de $\omega$ \unsure{impliquant} un iso canonique}
\margin{$\mathbb Z/m\mathbb Z\overset\sim\longrightarrow D^+$}
Les \deleted{points de} sommets du polygône seront désignés indifféremment par
$\widetilde R_i$ $(i\in\widetilde{\unsure{I|J}})$.
On n'a pas utilisé encore le fait que $m$ est pair, ce qui équivaut au fait que
l'ens des $\widetilde R_i$ dans $V$ est stable par la symétrie
$x\mapsto\deleted{\unsure{+}}-x$
